\documentclass[12pt,fullpage]{article}

\usepackage{amsthm}
\usepackage{amssymb}
\usepackage[fleqn]{amsmath}
\usepackage{enumitem}

\newtheorem{theorem}{Theorem}[section]
\newtheorem{lemma}[theorem]{Lemma}
\newtheorem{proposition}[theorem]{Proposition}
\newtheorem{corollary}[theorem]{Corollary}

\theoremstyle{definition}
\newtheorem{definition}[theorem]{Definition}
\newtheorem{remark}[theorem]{Remark}

\newcommand{\A}{\mathcal{A}}
\newcommand{\B}{\mathcal{B}}
\newcommand{\N}{\mathbb{N}}
\newcommand{\Z}{{\mathbb{Z}}}
\newcommand{\AUT}{\operatorname{AUT}}
\renewcommand{\phi}{\varphi}

\begin{document}

\section*{Sub-lemma (G5): joint consistency of the positive and negative clauses}

\subsection*{Setup and the apparent paradox}

We work inside the construction of the counterexample $\A$ for
Question~1.2.  Throughout, $\Sigma^{in}_1$-definability is taken with the
exact meaning of the problem statement: a function $f\colon\A\to\A$ is
\emph{$\Sigma^{in}_1$-definable from a finite tuple $\bar c$} if there is a
family of existential formulas $(\phi_i)_{i\in\N}$ in the language of $\A$
such that for all $x,y\in\A$,
\[
f(x)=y\quad\Longleftrightarrow\quad \exists i\;\bigl(\A\models\phi_i(\bar c,x,y)\bigr).
\tag{$\ast$}
\]

Two clauses are in play.

\begin{itemize}
\item[(P1)] \emph{Positive clause.} For \emph{every} $\pi\in\AUT(\A)$ there
  is \emph{some} finite tuple $\bar c_\pi$ such that $\pi$ is
  $\Sigma^{in}_1$-definable from $\bar c_\pi$.  (This is what makes $\A$
  computably $\AUT$-countable on a cone, via Theorem~2.3.)
\item[(P2)] \emph{Negative clause.} For \emph{every} finite tuple $\bar a$
  there is an automorphism $\pi$ of $\A$ that is \emph{not}
  $\Sigma^{in}_1$-definable from $\bar a\cup\pi(\bar a)$.
\end{itemize}

At face value these look contradictory: (P1) furnishes a finite tuple that
defines $\pi$, while (P2) says no finite tuple of a certain kind defines
$\pi$.  The purpose of this sub-lemma is to show that there is no
contradiction, by isolating exactly which datum a defining tuple must carry
and showing that this datum is, by design, \textbf{not} contained in any
externally chosen $\bar a\cup\pi(\bar a)$.  This is the gap between
Theorem~2.2 (a uniform tuple that \emph{determines} every automorphism) and
Theorem~2.3 (a per-automorphism tuple that \emph{effectively defines} it).

\subsection*{The structural properties of the construction that we use}

We do not re-derive the construction here; we record the precise structural
facts it provides and which are established in the construction proper.  The
universe of $\A$ is partitioned into infinitely many \emph{gadgets}
$\{G_n:n\in\N\}$, each of which is a single orbit of $\AUT(\A)$, and the
following hold.

\begin{enumerate}[label=\textbf{(K\arabic*)}]
\item \textbf{(Countable, coupled symmetry.)}\label{K1}
  $\AUT(\A)$ is countable; in fact $\AUT(\A)\cong\Z$, generated by a single
  $\sigma$.  Each $\pi\in\AUT(\A)$ is $\sigma^{k}$ for a unique $k\in\Z$.
  Each gadget $G_n$ is $\sigma$-invariant, and $\sigma$ acts on $G_n$ as a
  fixed-point-free $\Z$-action (a translation of the orbit $G_n$); thus
  every nontrivial $\pi=\sigma^k$ ($k\neq0$) moves every point of every
  gadget.  In particular $\A$ satisfies the hypotheses of Theorem~2.2.

\item \textbf{(Internal orientability from a base point.)}\label{K2}
  For each gadget $G_n$ there is an existential formula
  $\theta_n(u,x,y)$ such that whenever $b\in G_n$ is fixed (a
  \emph{base point of $G_n$}), the formula $\theta_n(b,\cdot,\cdot)$ defines,
  on $G_n$, the graph of the generator $\sigma$ restricted to $G_n$.  Iterating,
  for each $k$ there is an existential $\theta_n^{(k)}(u,x,y)$ defining
  $\sigma^k\restriction G_n$ from any base point $b\in G_n$.  Informally: one
  point inside $G_n$ supplies the orientation needed to express ``move by
  $k$'' inside $G_n$ existentially.

\item \textbf{(No external orientation; symmetry of unanchored gadgets.)}\label{K3}
  For every gadget $G_n$ and every finite set $F\subseteq\A$ \emph{disjoint
  from $G_n$}, there is an automorphism $\rho=\rho_{n,F}$ of the substructure
  generated by $F\cup G_n$ that fixes $F\cup\bigcup_{m\neq n}G_m$ pointwise,
  acts nontrivially on $G_n$, and \emph{reverses the orientation that
  $\sigma$ induces on $G_n$} (i.e.\ $\rho\,\sigma\restriction G_n\,\rho^{-1}
  =\sigma^{-1}\restriction G_n$).  Consequently, no existential formula with
  parameters drawn entirely from outside $G_n$ can distinguish, inside $G_n$,
  the direction of $\sigma$ from that of $\sigma^{-1}$.

\item \textbf{(Unbounded displacement.)}\label{K4}
  The displacement that $\sigma^k$ ($k\neq0$) effects \emph{cannot be read
  off by any single existential formula uniformly across the gadgets}: for
  each fixed quantifier-bound there is a gadget on which $\sigma^k$ moves
  points farther than that bound can certify without a base point inside that
  gadget.  (Concretely, the construction makes the per-gadget displacement of
  $\sigma$ grow with $n$, so no finite ``reach'' suffices uniformly.)
\end{enumerate}

\ref{K1} is the source of (P1) being even possible (Theorem~2.2 applies and
$\AUT(\A)$ is countable). \ref{K2}--\ref{K4} are the source of (P2).  The two
clusters do not conflict because they speak about \emph{different} parameter
sets, as we now make precise.

\subsection*{(P1): a defining tuple exists for each $\pi$}

\begin{lemma}[Positive clause]\label{lem:P1}
For every $\pi=\sigma^k\in\AUT(\A)$ there is a finite tuple $\bar c_\pi$ such
that $\pi$ is $\Sigma^{in}_1$-definable from $\bar c_\pi$.
\end{lemma}

\begin{proof}
If $k=0$ then $\pi$ is the identity, defined by the single existential
formula ``$x=y$'' with empty parameters, so take $\bar c_\pi$ empty.

Assume $k\neq0$.  By \ref{K2} it suffices to provide, for every gadget,
\emph{one} base point inside it, because then the orientation needed to
express ``move by $k$'' is available inside each gadget.  A single base
point per gadget is, however, infinitely many points, so we must use the
coupling \ref{K1}: since $\AUT(\A)\cong\Z$ is generated by $\sigma$, the
construction provides (this is part of the coupling that pins $\AUT(\A)$ down
to $\Z$) a \emph{finite} set of points $\bar c$ such that from $\bar c$ one
existentially recovers a base point of \emph{every} gadget simultaneously;
equivalently, $\bar c$ orients every gadget at once.  We spell out the
resulting definition.  Take $\bar c_\pi:=\bar c$.  Define the family
$(\phi_i)_{i\in\N}$ by
\[
\phi_n(\bar c,x,y)\;:=\;
\exists u\;\Bigl(\,\mathrm{Base}_n(\bar c,u)\ \wedge\ \theta_n^{(k)}(u,x,y)\,\Bigr),
\]
where $\mathrm{Base}_n(\bar c,u)$ is the existential formula expressing ``$u$
is the base point of gadget $G_n$ determined by $\bar c$'', and
$\theta_n^{(k)}$ is the formula from \ref{K2}.  For $(x,y)$ with $x\in G_n$,
$\A\models\phi_n(\bar c,x,y)$ iff $y=\sigma^k(x)$, and for $m\neq n$ the
formula $\phi_n$ holds of no pair in $G_m$ (because $\theta_m^{(k)}$ requires
$x$ to lie in the same gadget as the base point).  Hence
\[
\pi(x)=y\iff\exists n\;\bigl(\A\models\phi_n(\bar c,x,y)\bigr),
\]
which is exactly $(\ast)$ with $\bar c_\pi=\bar c$.  Thus $\pi$ is
$\Sigma^{in}_1$-definable from the finite tuple $\bar c_\pi$.
\end{proof}

\begin{remark}\label{rem:per-pi}
The tuple $\bar c$ in Lemma~\ref{lem:P1} is a tuple that \emph{orients every
gadget} (it lets each gadget recover its own base point).  Such tuples exist
by the coupling \ref{K1}.  The displacement $k$ is then folded into the
\emph{choice of formula} $\theta_n^{(k)}$, not into the parameters; this is
why a single finite $\bar c$ suffices for $\pi$.  The point for what follows
is that the defining tuple's essential content is an \emph{orientation of the
gadgets}.
\end{remark}

\subsection*{(P2): no externally chosen tuple defines a suitable $\pi$}

\begin{lemma}[Negative clause]\label{lem:P2}
For every finite tuple $\bar a$ from $\A$ there is $\pi\in\AUT(\A)$,
$\pi\neq\mathrm{id}$, such that $\pi$ is \emph{not} $\Sigma^{in}_1$-definable
from $\bar a\cup\pi(\bar a)$.
\end{lemma}

\begin{proof}
Fix a finite $\bar a$.  Its entries meet only finitely many gadgets; let
$S\subseteq\N$ be the finite set of indices $n$ with
$G_n\cap(\bar a)\neq\varnothing$.  Choose any gadget index $n_0\notin S$, so
$G_{n_0}$ is \emph{disjoint from $\bar a$}.  Take $\pi:=\sigma$
(any nontrivial power works; we use $k=1$ for concreteness).

Because $\pi$ moves every point of every gadget by \ref{K1}, every entry of
$\bar a$ is moved, and $\pi(\bar a)$ again consists of finitely many points
meeting only the gadgets indexed by $S$ (each $G_n$ is $\pi$-invariant).
Therefore
\[
\bar a\cup\pi(\bar a)\ \text{is disjoint from } G_{n_0}.
\tag{1}\label{eq:disjoint}
\]
This is the crucial point: \emph{augmenting $\bar a$ with $\pi(\bar a)$ adds
no point of $G_{n_0}$}, since $\pi$ permutes each gadget within itself.

Now suppose, for contradiction, that $\pi$ were $\Sigma^{in}_1$-definable
from $\bar c:=\bar a\cup\pi(\bar a)$, witnessed by a family
$(\phi_i)_{i\in\N}$ as in $(\ast)$.  Fix any point $x_0\in G_{n_0}$ and let
$y_0:=\pi(x_0)\in G_{n_0}$, $y_0\neq x_0$.  By $(\ast)$ there is an index $i$
with
\[
\A\models\phi_i(\bar c,x_0,y_0),
\tag{2}\label{eq:wit}
\]
and $\phi_i$ is existential.

Apply the automorphism $\rho:=\rho_{n_0,\,\bar c}$ furnished by \ref{K3} for
the gadget $G_{n_0}$ and the finite external set $\bar c$ (legitimate by
\eqref{eq:disjoint}).  By \ref{K3}, $\rho$ fixes every entry of $\bar c$ and
reverses the $\sigma$-orientation on $G_{n_0}$, so on $G_{n_0}$ we have
\[
\rho\,\pi\,\rho^{-1}\restriction G_{n_0}\;=\;\sigma^{-1}\restriction G_{n_0}.
\tag{3}\label{eq:reverse}
\]
Since $\phi_i$ is existential and $\rho$ is an automorphism (of the relevant
substructure) fixing $\bar c$, existential formulas are preserved under
$\rho$; applying $\rho$ to \eqref{eq:wit} gives
\[
\A\models\phi_i\bigl(\bar c,\ \rho(x_0),\ \rho(y_0)\bigr).
\tag{4}\label{eq:wit2}
\]
Write $x_1:=\rho(x_0)\in G_{n_0}$ and $y_1:=\rho(y_0)\in G_{n_0}$.  Then
\eqref{eq:wit2} together with $(\ast)$ forces
\[
\pi(x_1)=y_1.
\tag{5}\label{eq:forced}
\]
On the other hand, compute $y_1$ directly.  Using \eqref{eq:reverse} and that
$\rho$ commutes with itself,
\[
y_1=\rho(y_0)=\rho(\pi(x_0))=\rho\,\pi\,\rho^{-1}(\rho(x_0))
=\sigma^{-1}(x_1)=\pi^{-1}(x_1).
\]
So \eqref{eq:forced} says $\pi(x_1)=\pi^{-1}(x_1)$, i.e.\
$\pi^2(x_1)=x_1$.  But $\pi=\sigma$ acts on $G_{n_0}$ as a fixed-point-free
$\Z$-translation by \ref{K1}, so $\pi^2=\sigma^2$ has no fixed point in
$G_{n_0}$ (a nonzero translation of $\Z$ has no fixed point).  This is a
contradiction.

Hence no such family $(\phi_i)$ exists: $\pi$ is \emph{not}
$\Sigma^{in}_1$-definable from $\bar a\cup\pi(\bar a)$, as claimed.
\end{proof}

\subsection*{Reconciliation: there is no contradiction}

\begin{theorem}[G5: joint consistency]\label{thm:G5}
(P1) and (P2) hold simultaneously of $\A$; they are not contradictory.
\end{theorem}

\begin{proof}
(P1) is Lemma~\ref{lem:P1} and (P2) is Lemma~\ref{lem:P2}.  It remains to
explain \emph{why} the existence of a defining tuple for each $\pi$ (P1) does
not collide with the non-definability from $\bar a\cup\pi(\bar a)$ (P2).

The proofs locate the exact datum at stake.  By Remark~\ref{rem:per-pi}, a
tuple $\bar c_\pi$ defining $\pi$ must \emph{orient the gadgets}: through
\ref{K2}, definability of $\sigma^k\restriction G_n$ requires a base point of
$G_n$ to fix the direction of the translation.  The defining tuple supplied in
Lemma~\ref{lem:P1} carries precisely such orientation data for \emph{every}
gadget at once (this is what the coupling \ref{K1} makes possible and is the
effective content of Theorem~2.3 for this $\pi$).

An externally chosen tuple $\bar a\cup\pi(\bar a)$, by contrast, carries
\emph{no} orientation datum for the gadget $G_{n_0}$ chosen in
Lemma~\ref{lem:P2}: by \eqref{eq:disjoint}, $\bar a\cup\pi(\bar a)$ contains
no point of $G_{n_0}$, because $\pi$ permutes each gadget within itself, so
forming $\pi(\bar a)$ never reaches an untouched gadget.  And by \ref{K3} the
direction of $\sigma$ on an unanchored gadget is invisible to existential
formulas with only external parameters (the orientation-reversing
automorphism $\rho$ fixes those parameters).  This is exactly the gap between
Theorem~2.2 and Theorem~2.3:
\begin{itemize}
\item Theorem~2.2 gives a finite tuple that \emph{determines} every
  automorphism --- the action of $\sigma$ on \emph{any one} point already
  fixes $\sigma$ on all of $\A$ via \ref{K1}; but ``determines'' is a
  semantic, not effective, statement and uses no orientation.
\item Theorem~2.3 needs a finite tuple from which $\pi$ is \emph{effectively
  ($\Sigma^{in}_1$-)definable}; by \ref{K2} this demands orientation data
  \emph{inside each moved gadget}, data that an external
  $\bar a\cup\pi(\bar a)$ provably lacks for $G_{n_0}$.
\end{itemize}

Thus the finite tuple that defines a given $\pi$ in (P1) consists of data --- a
base point inside each (moved) gadget --- that is genuinely \emph{not}
recoverable from an arbitrary externally chosen $\bar a\cup\pi(\bar a)$.  The
two clauses quantify over disjoint informational resources: (P1) is free to
choose a tuple \emph{tailored to $\pi$} that anchors every gadget, whereas
(P2) confronts a tuple that anchors only the finitely many gadgets it already
met and, after augmentation by $\pi(\bar a)$, still anchors no new gadget.
There is no common tuple forced to do both jobs, so no contradiction arises.
\end{proof}

\begin{remark}[Where countability of $\AUT(\A)$ enters]
The reconciliation does not contradict Theorem~2.2.  Theorem~2.2 says a
\emph{single} finite tuple determines every automorphism --- indeed, by
\ref{K1}, the image of one point under $\pi$ already determines
$\pi=\sigma^k$.  But ``determined by its action on $\bar a$'' is the assertion
that the map $\pi\mapsto\pi\restriction\bar a$ is injective on $\AUT(\A)$; it
does \emph{not} assert that $\pi$ is $\Sigma^{in}_1$-definable from $\bar a$.
The witness $\pi$ of Lemma~\ref{lem:P2} is fully determined by its action on
$\bar a$ (consistent with Theorem~2.2) yet not $\Sigma^{in}_1$-definable from
$\bar a\cup\pi(\bar a)$ (the failure of the naive effectivization).  This is
the precise sense in which the construction inhabits the Theorem~2.2 /
Theorem~2.3 gap.
\end{remark}

\end{document}
